\section{Conclusion}

MonopolyBench makes strategic execution and social interaction inspectable at the decision where they occur. The twenty-game cohort shows that recipients often decline exchanges completing their own sets, legal financing can outlast an apparent cash shortfall, and later opportunities depend on construction, transfers and bankruptcy. It also shows why these observations need careful attribution: defaults can resemble restraint, submitted plans can disagree with received actions, and acceptance can accompany explicit rejection of a proposal's premise.

We connect complete offer and transfer populations, source-linked decision text and historical legal alternatives. All 50,918 original events replay; the portable suite checks sixty-six defaults and sixty eliminations. The separate reference-policy study reproduces its own twenty streams while exposing scenario and comparator weaknesses that must be addressed before a broader skill evaluation. Reliable execution of an experiment is necessary but does not validate its strategic target.

The eight research questions therefore have different kinds of answers. This study establishes local feasibility, observed tradeoffs, plan and response records, repeated errors and collective transfers. It does not establish optimal play, latent beliefs, causal persuasion, durable learning or generalization to unfamiliar opponents. The broader relationship between general and multi-agent capability remains unmeasured. Stating those boundaries alongside the positive results gives subsequent experiments explicit targets, rather than asking a win count or a compelling transcript to answer every question.

